%%BEGIN_ABSTRACT
Latent world models are commonly evaluated with a state probe, a regressor fitted on encoded observations and then applied to predicted latents. We show that norm drift distorts this evaluation. During autoregressive prediction, the norm of a predicted latent can drift while its direction stays accurate, and probes fitted on encodings are highly sensitive to the norm. In synthetic latents that store the state in their direction, a norm error of $5$\% per prediction step compounds within a few dozen steps until the probe's error is $27$ times that of always predicting the mean state. When the norm varies independently of the true error, the probe orders only about $61$\% of pairs of latents correctly, where chance is 50\%. Of $39$ evaluations of predicted latents that we surveyed, only one reports prediction norms. The same distortion appears in a published model. For TCRL on pendulum, rescaling the predictions to the norm of the encodings removes $97$\% of the probe's error at the end of the training horizon, because the predicted directions are accurate and only their norms are off. When the norm dominates the probe score, a worse model that drifts less can rank above a better one. We propose re-projection, which rescales each prediction to a reference norm taken from the encodings, keeps its direction and applies the same probe again. We add four checks that test whether the re-projected error reflects what the model predicts. Re-projection needs no retraining and can be added to any state-probe evaluation of predicted latents.
%%END_ABSTRACT
